% \section{Robot Architecture}
% FRIDA's software runs on ROS~2 Jazzy and is split into four modules (vision, HRI, manipulation and navigation), each in its own Docker container. A launcher detects whether it runs on a Jetson, a CUDA workstation or a CPU-only machine and starts the matching images, which CI builds on every push and pull request and publishes to Docker Hub. Modules communicate over CycloneDDS; speech, TTS and the LLM servers run outside ROS and are reached over gRPC or HTTP. Each task is coordinated by a task manager: most are finite state machines, while GPSR executes behavior trees built at run time from the interpreted commands. Task managers reach the modules only through subtask managers that wrap each module's ROS~2 interfaces. Every call waits for its service with a timeout and returns a status code instead of raising, so an unavailable module becomes a failure the task can handle (retry, skip or continue) rather than a crash; the same layer lets any module be mocked in tests.
\section{Robot Architecture}

\begin{figure}[H]
\centering
\begin{tikzpicture}[font=\scriptsize, node distance=3mm and 2mm,
  box/.style={draw, rounded corners=1pt, minimum height=5mm, align=center, inner sep=2pt},
  mod/.style={box, minimum width=20mm, fill=gray!10}]
\node[box, minimum width=86mm, fill=gray!25] (tm) {Task managers: FSM per task \;---\; GPSR: behavior tree built at run time};
\node[box, minimum width=86mm, below=of tm] (sm) {Subtask managers: timeout-guarded, status codes, \texttt{@mockable}};
\node[mod, below=of sm, xshift=-33mm] (vis) {Vision};
\node[mod, right=of vis] (hri) {HRI};
\node[mod, right=of hri] (man) {Manipulation};
\node[mod, right=of man] (nav) {Navigation};
\node[box, below=of $(vis.south)!0.5!(hri.south)$, minimum width=40mm, text width=38mm, dashed] (grpc) {STT, TTS, LLM servers\\(gRPC / HTTP, outside ROS)};
\node[box, below=of $(man.south)!0.5!(nav.south)$, minimum width=40mm, text width=38mm] (hw) {Drivers: arm, gripper,\\cameras, LiDARs, base};
\draw[->] (tm) -- (sm);
\foreach \m in {vis,hri,man,nav} \draw[->] (sm.south -| \m.north) -- (\m.north);
\draw[->] (hri) -- (grpc.north -| hri.south);
\draw[->] (man.south) -- (hw.north -| man.south);
\draw[->] (nav.south) -- (hw.north -| nav.south);
\end{tikzpicture}
\caption{Software architecture. Each module runs in its own container and communicates over CycloneDDS with shared memory.}
\label{fig:arch}
\end{figure}

FRIDA's software runs on ROS~2 Jazzy (JetPack~7 on two NVIDIA Jetson AGX Orin) and is split into four modules (vision, HRI, manipulation and navigation), each packaged in its own Docker container (Fig.~\ref{fig:arch}).
\subsubsection{Deployment.} A single launcher detects whether it runs on a Jetson, a CUDA workstation or a CPU-only machine and starts the matching images, so the same commands run on the robot and on development laptops. Containers are pinned to disjoint CPU sets (navigation on cores 0--3, perception and planning on 4--11) so that heavy perception workloads cannot starve the navigation control loop.

\subsubsection{Communication.} Modules communicate over CycloneDDS; on the robot, shared-memory transport avoids copying camera images and point clouds between containers. Speech, TTS and LLM servers run outside ROS and are reached over gRPC or HTTP, keeping their Python and CUDA dependencies isolated from the ROS images.

\subsubsection{Task execution.} Each task is coordinated by a task manager: structured tests use finite state machines, while GPSR executes behavior trees built at run time from the interpreted commands (Sec.~\ref{sec:scheduling}). Task managers reach the modules only through subtask managers that wrap each module's ROS~2 interfaces. Every call waits for its service with a timeout and returns a status code instead of raising, so an unavailable module becomes a failure the task can handle (retry, skip or continue) rather than a crash.

% \subsubsection{Testing.} Every subtask function is declared mockable, and an automated test enforces it, so any module can be replaced by mocks to test a task manager without the robot. CI builds the ROS~2 workspace on every push and pull request, and rebuilds and publishes the images to Docker Hub whenever their Dockerfiles change.